\scanpage{121}
\mspage{- 93 -}
n\textsuperscript{o} 371

% L1–2
$\handwritten{\bigwedge}(N)$, où $N$ est un sous-module de type fini de $\unsure{M}$ ; il suffit alors d'appliquer le cor.3 de la prop.2 du n\textsuperscript{o}1.

% L3–6
COROLLAIRE 5.- \uline{Soient} $N$, $P$ \uline{deux sous-modules d'un} $A$-\uline{module} $M$ \uline{tels que} $N\cap P$ \uline{soit facteur direct de} $N$ \uline{et de} $P$ \uline{et} $N+P$ \uline{facteur direct de} $M$. \uline{Alors} $\handwritten{\bigwedge}(N)$ \uline{et} $\handwritten{\bigwedge}(P)$ \uline{s'identifient canoniquement à des sous-algèbres de} $\handwritten{\bigwedge}(M)$, \uline{et} $\handwritten{\bigwedge}(N\cap P)$ \uline{à l'intersection} $\handwritten{\bigwedge}(N)\cap \handwritten{\bigwedge}(P)$.

% L7–12
En effet, si $E$, $F$ sont des supplémentaires de $N\cap P$ dans $N$ et $P$ respectivement et $G$ un supplémentaire de $N+P$ dans $M$, $M$ est somme directe des quatre sous-modules $N\cap P$, $E$, $F$, $G$ ; donc $\handwritten{\bigwedge}(N)$, $\handwritten{\bigwedge}(P)$ et $\handwritten{\bigwedge}(N\cap P)$ s'identifient à des sous-algèbres de $\handwritten{\bigwedge}(M)$ (cor.3). En outre, en vertu du cor.1, $\overset{n}{\handwritten{\bigwedge}}\,N$ (resp. $\overset{n}{\handwritten{\bigwedge}}\,P$) s'identifie à la somme directe des sous-modules $\overset{p}{\handwritten{\bigwedge}}(N\cap P)\otimes\overset{n-p}{\handwritten{\bigwedge}}\,E$ (resp. $\overset{p}{\handwritten{\bigwedge}}(N\cap P)\otimes\overset{n-p}{\handwritten{\bigwedge}}\,F$) pour $0\leq p\leq n$ ; on a donc bien $\handwritten{\bigwedge}(N\cap P)=\handwritten{\bigwedge}(N)\cap \handwritten{\bigwedge}(P)$.

% L13–14
PROPOSITION 8.\unsure{=|-} \uline{Soit} $(e_{\lambda})_{\lambda\in L}$ \uline{une base d'un} $A$-\uline{module} $M$ ; \uline{munissons} $L$ \uline{d'une structure d'ordre total}, \uline{et pour toute partie finie} $H$ \uline{de} $L$, \uline{posons}
% L15
\begingroup\csname tagsleft@true\endcsname
\[
 e_H=e_{\lambda_1}\handwritten{\wedge} e_{\lambda_2}\handwritten{\wedge}\mathord{..}\handwritten{\wedge} e_{\lambda_n}\tag*{(11)}
\]
\endgroup
% L16–18
\uline{si} $(\lambda_i)_{1\leq i\leq n}$ \uline{est la suite strictement croissante des éléments de} $H$. \uline{Alors les éléments homogènes} $e_H$, \uline{où} $H$ \uline{parcourt l'ensemble} $F(L)$ \uline{des parties finies de} $L$, \uline{forment une base de} $\handwritten{\bigwedge}(M)$.

% L19–23
Supposons d'abord $L$ \uline{fini}, soit par exemple $L=\unsure{[1,m]|[\![1,m]\!]}$ muni de l'ordre usuel ; alors la proposition résulte de la prop.7, car si on pose $M_i=Ae_i$ pour $1\leq i\leq m$, $M$ est somme directe des $M_i$ et il résulte aussitôt du cor.3 de la prop.2 du n\textsuperscript{o}1 que $\handwritten{\bigwedge}(M_i)=A\oplus M_i$ (autrement dit $\overset{p}{\handwritten{\bigwedge}}M_i=0$ pour $p\geq 2$) ; la conclusion résulte alors du cor.1 de la prop.7 et du chap.II, 3e éd., \S~3, n\textsuperscript{o}4, prop.4.

% L24–25
Supposons ensuite $L$ infini et pour toute partie finie $J$ de $L$, soit $M_J$ le sous-module de $M$ ayant pour base les $e_\lambda$ tels que $\lambda\in J$ ; il est clair que $M$ est limite
% archive sheet number: 120
